\uniterablesection{Висновки}

За результатами проходження передкваліфікаційної практики на базі кафедри інженерії програмного забезпечення та штучного інтелекту НУВГП проаналізовано предметну область, сформульовано вимоги, формалізовано архітектуру й алгоритми вебзастосунку для тривимірного моделювання внутрішньобудинкової електромережі приміщень та перевірено коректність його розрахунків. Отримано такі результати.

\begin{enumerate}[noitemsep, topsep=0.3em, leftmargin=*, label=\arabic*.]
    \item Проаналізовано проблематику планування розміщення електрообладнання, трасування проводки та розрахунку оздоблювальних матеріалів. Показано, що довжина проводу визначається ортогональним маршрутом, а не прямою відстанню: на контрольній сцені маршрут через стелю (7,10~м) більш ніж удвічі довший за пряму (близько 3,00~м). Задачу розкрою листового матеріалу формалізовано як покриття чистої площі поверхонь гільйотинними розрізами з повторним використанням відрізків і виведено нижню межу кількості листів за площею.

    \item Виконано огляд трьох редакторів планування приміщень (Sweet Home 3D, Planner 5D, RoomSketcher) і двох онлайн-калькуляторів гіпсокартону. Встановлено, що жоден із них не поєднує тривимірну модель, трасування проводів із розрахованою довжиною та розкладку листів; на цій підставі визначено нішу застосунку.

    \item Розроблено специфікацію вимог за ISO/IEC/IEEE 29148:2018: чотири групи зацікавлених сторін, 14~функціональних і 8~нефункціональних вимог із критеріями приймання (структуровано за ISO/IEC~25010:2023), діаграму варіантів використання та опис ключового сценарію.

    \item Обґрунтовано архітектуру клієнтського застосунку без збірки, сервера та бази даних (JavaScript, ES-модулі, Three.js); описано діаграму C4, модель даних із форматом JSON версії~2, нормалізацією під час імпорту та скасуванням на 60 кроків. Алгоритм розкрою винесено в окремий модуль, що не залежить від DOM і графічної бібліотеки.

    \item Формалізовано геометричну модель приміщення зі стіновою системою координат, алгоритм ортогонального трасування проводів і формули розрахунків площ, листів ГКЛ, профілів каркаса стелі, плінтусів і довжини проводки. Описано алгоритм розкрою з пулом відрізків, гільйотинним розрізом, складанням зі шматків і режимом «Економія».

    \item Верифіковано розрахунки: на контрольній сцені з двох кімнат усі 16 показників застосунку збіглися з незалежним перерахунком окремим скриптом. Режими «Горизонтально» та «Економія» досягли нижньої межі в 30~листів проти 33 за простою оцінкою (економія 9~\%); час розрахунку розкладки склав від 0,05 до 68~мс.

    \item Визначено обмеження поточної версії: відсутні електротехнічні розрахунки та автоматизовані тести, розкрій є евристичним, назви режимів розкладки в інтерфейсі потребують уточнення, імпорт недовірених файлів потребує екранування тексту, Three.js залежить від мережевого CDN.

    \item Опрацьовано 26 джерел: рецензовані публікації з теорії розкрою, нормативний документ, стандарти інженерії програмного забезпечення, технічну документацію та описи програмних аналогів. Записи наукових публікацій перевірено за реєстром Crossref, а веб-джерела~-- за сторінками виробників.
\end{enumerate}

\vspace{0.5em}
Отримані результати утворюють методологічну й проєктну основу для кваліфікаційної роботи. Наступними етапами є: реалізація розрахункового ядра електротехнічних параметрів (розрахункове навантаження, переріз жил, падіння напруги, струми короткого замикання, вибір апаратів захисту) із прив'язкою до довжин проводів, які вже визначає застосунок; оформлення розрахунків як автоматизованих тестів; локальне підключення Three.js для офлайн-роботи; усунення виявлених обмежень безпеки та інтерфейсу; юзабіліті-тестування за вимогою NFR-08.
